\scanpage{27}
\mspage{25}

A titre accessoire il peut être utile (peut-être) que $i$ soit une ``cofibration''$^{+}$. On définit $\Omega(S)$ comme la fibre de $\mathrm{Ch}(S)$ sur $S\times S$ (par rapport à l'unique morphisme $e=e_M\mathbin{\unsure{-|\longrightarrow}}S\times S$, où $e_M$ est l'objet \deleted{\unsure{neutre}} \inserted{nul} de $M$).

% slanted
\margin{et peut-être que $S\longrightarrow e$ \unsure{soit|doit} également une ``fibration'' ?}

\diagram{\begin{tikzcd}[ampersand replacement=\&,row sep=small]
 \mathrm{Ch}(S)\arrow[d]\&\Omega S\arrow[l]\arrow[d]\\
 S\times S\&e\arrow[l]\quad .
\end{tikzcd}}

C'est un objet de $M$, qu'on regardera surtout comme un objet de $MW^{-1}$. Il faut s'assurer qu'on a une \deleted{\unsure{isom.}} \inserted{bij.} fonctorielle en $T$ dans $MW^{-1}$ :

\diagram{\begin{tikzcd}[ampersand replacement=\&,row sep=small,column sep=tiny]
 \unsure{\Pi|\pi}_1\underline{\mathrm{Ch}}_{\infty}^{W}(T,S)\arrow[d,equal,"\text{déf}"]\&[-5pt]\simeq\&
 \mathrm{Hom}_{MW^{-1}}(T,\Omega S)\arrow[d,equal,"\text{déf}"]\\
 \Pi_1(T,S)\&\&\Pi_0(\underline{\mathrm{Ch}}^{W}_{\infty}(T,\Omega S)).
\end{tikzcd}}

Pour cela il s'impose de comparer

\diagram{\begin{tikzcd}[ampersand replacement=\&,row sep=large,column sep=tiny]
 \underline{\mathrm{Ch}}^{W}_{\infty}(T,\Omega S)\&\text{et}\&
 \Omega_{\mathrm{Hot}_{\bullet}}\bigl(\underline{\mathrm{Ch}}^{W}_{\infty}(T,S)\bigr)
 \arrow[d,equal,"{\substack{\text{``par ``définition''}\\\text{pour }\Omega\text{ de }\mathrm{Hot}_{\bullet}}}"]\\
 \&\&\underline{\mathrm{Ch}}_{\infty}\bigl(\underline{\mathrm{ch}}^{W}_{\infty}(T,S);\unsure{\sigma|0|o},\sigma\bigr),
\end{tikzcd}}

puisqu'il s'agit \deleted{\unsure{de}} \unsure{d'y|de} construire une bijection entre leurs $\Pi_0$.
% archive sheet number: 26
